\clearpage\begingroup\SourceLeaf{587}
\ArithmeticHead{32}{L}{PHILIPP MAENNCHEN, GAUSS ALS ZAHLENRECHNER.}
\renewcommand{\thefootnote}{\arabic{footnote})}\setcounter{footnote}{0}
\begin{center}VII.\\[3pt]{\large\textls[80]{Die Tafel des quadratischen Charakters der Primzahlen.}}\end{center}
\noindent\hspace*{\parindent}Bei dem Bestreben, grosse Zahlen in Faktoren zu zerlegen, ist der im
vorigen Abschnitt gekennzeichnete Weg nicht der erste und nicht der einzige
gewesen. Die frühesten Versuche bestanden wohl darin, dass er Zahlengruppen
von der Form $a^2+1$, $a^2+2$, $a^2+3$, $a^2+4$, $a^2+5$, $a^2+7$, $a^2+9$, $a^2+11$,
allgemein $a^2+p$ bezw. $a^2+p^2$ aufstellte und alle Zahlen dieser Gruppen in
ihre Primfaktoren zerlegte. Die Beobachtungen, die er bei dieser Gelegenheit machte, führten ihn sowohl zu der jetzt zu besprechenden Theorie der
quadratischen Reste, als auch zu der \textit{Tafel zur Cyklotechnie}, von der im
nächsten Abschnitt die Rede sein wird.

Die Beobachtungen über quadratische Reste und Nichtreste waren zwar
schon vor \textsc{Gauss} von \textsc{Fermat}, \textsc{Euler} und \textsc{Legendre} gemacht worden, aber
dem jugendlichen Forscher stand zu jener Zeit noch keine Literatur zur Verfügung\footnote{\textsc{Euler}s wichtigste Arbeiten auf diesem Gebiet:
\begin{list}{\arabic{enumi}.}{\usecounter{enumi}\setlength{\leftmargin}{1.5em}\setlength{\itemsep}{0pt}\setlength{\topsep}{1pt}}
\item \textit{Theoremata circa divisores numerorum in hac forma $paa+qbb$ contentorum.} Comm. acad. Petrop. \textbf{14} (1744/6), 1757, S. 151—181; \textsc{Leonhardi Euleri}, Opera omnia, series I, vol. II, S. 194—222.
\item \textit{De numeris qui sunt aggregata duorum quadratorum.} Novi comm. acad. Petrop. \textbf{4} (1752/3), 1758, S. 3 bis 40; Opera, ibid., S. 295—327.
\item \textit{Specimen de usu observationum in mathesi pura.} Novi comm. acad. Petrop. \textbf{6} (1756/7), 1761, S. 185 bis 230; Opera, ibid., S. 459—492.
\item \textit{Observationes circa divisionem quadratorum per numeros primos.} Opuscula analytica I, 1783, S. 64—84; Opera, series I, vol. III, S. 497.
\end{list}}.

\textsc{Gauss} bemerkte wohl bald, dass zu jeder Zahl bestimmte Reste und Nichtreste gehören, dass das Produkt zweier Reste, sowie auch das zweier Nichtreste ein Rest, dagegen das Produkt eines Restes und eines Nichtrestes ein
Nichtrest ist, dass ferner die Untersuchung für eine zusammengesetzte Zahl
als Modul dieselben Reste und Nichtreste liefert, wie die für ihren kleinsten
Primteiler. Um nun einerseits diese Beobachtungen praktisch zu verwerten
und andererseits möglicher Weise noch tiefer in das neue Gebiet hineinschauen
zu können, unternahm er eine Arbeit, die wieder höchst charakteristisch für
seine Gewandtheit und Ausdauer ist. Er stellte eine Tabelle auf, die die
Primzahlen von 2 bis 997 als Reste in bezug auf die Primzahlen von 3 bis
\clearpage\SourceLeaf{588}\ArithmeticHead{33}{R}{VII. DIE TAFEL DES QUADRATISCHEN CHARAKTERS DER PRIMZAHLEN.}
\setcounter{footnote}{0}
\noindent 503 als Teiler enthält. Um einen Begriff von der Grösse dieses Unternehmens
zu geben, wollen wir nur durch einfache Abzählung feststellen, dass \textls[120]{16320mal
untersucht werden musste, ob eine Zahl quadratischer Rest oder
Nichtrest} war.\linebreak Entspricht nun der erzielte Gewinn dieser gewaltigen Arbeit? Wir dürfen diese Frage wohl bejahen. Denn abgesehen von der erreichten Gewandtheit in der Gewinnung von quadratischen Resten war \textsc{Gauss}
nun in der Lage, seine Methode der Exkludenten sehr weittragend auszugestalten, und wir wollen nachher an einem Beispiel zeigen, wie eine Faktorenzerlegung mit Hülfe dieser Tafel vor sich geht. Ausserdem aber wurde seine
Mühe glänzend belohnt durch den erzielten zahlentheoretischen Gewinn, nämlich durch die Erkenntnis des berühmten Satzes, den er »Theorema fundamentale in doctrina de residuis quadraticis« nannte. Dieses Gesetz, das Reziprozitätsgesetz der quadratischen Reste nach \textsc{Legendre}, von diesem und vor
ihm auch von \textsc{Euler} durch Induktion gefunden, wurde nun von \textsc{Gauss} zum
ersten Male bewiesen. \textsc{Gauss} gibt im ganzen acht verschiedene Beweise für
diesen Satz, woraus wir schliessen dürfen, welche grosse Bedeutung er ihm
beimisst\footnote{Vergl. P. \textsc{Bachmann}, \textit{Über Gauss’ zahlentheoretische Arbeiten}, Abhandlung 1 dieses Bandes, S. 14 ff. und 45 ff.}.

Doch nun zu der praktischen Verwendung für die Zahlenzerlegung! Wir
wählen das Beispiel, das \textsc{Gauss} in seinen \textit{Disquisitiones} wiederholt benutzt,
997331, d. i. $314159265:(9\mathord{.}5\mathord{.}7)$. Für diese Zahl ermittelt er die quadratischen Reste $-6,\;+13,\;-14,\;+17,\;+37,\;-53$. Auf Grund der Tafel
findet er, dass von allen Primzahlen bis 997 nur 127 die für 997331 gefundenen quadratischen Reste hat, dass also keine andere dieser Primzahlen
als 127 in 997331 ohne Rest enthalten sein kann.

In einem Brief an \textsc{Zimmermann}\footnote{Werke X, 1, S. 20.} spricht \textsc{Gauss} über dieses Verfahren und
bemerkt dabei »dies ist nämlich nur Eine einzelne Anwendung der Tafel«. Er
hat also anscheinend noch mancherlei andere Anwendungen im Auge gehabt;
allein im Nachlass hat sich nichts darüber gefunden. Doch gehen wir wohl
nicht fehl, wenn wir annehmen, dass er mit Hilfe dieser Tafel seine Methode
der Exkludenten zu höchster Leistungsfähigkeit zu steigern vermochte, so dass
er sich dieses Hilfsmittels bedienen konnte bei der Lösung der Aufgabe
\par\noindent X 2 Abh. 6.\hfill 5
\clearpage\SourceLeaf{589}\ArithmeticHead{34}{L}{PHILIPP MAENNCHEN, GAUSS ALS ZAHLENRECHNER.}
\setcounter{footnote}{0}
\noindent $n=ax^2+by^2$, bei der \textsc{Pell}schen Gleichung, sowie bei einer Aufgabe, die
uns im nächsten Abschnitt entgegentreten wird. Die Tafel vermag auch oft
sehr einfach darüber Aufschluss zu geben, dass eine Zahl zusammengesetzt
ist. Z. B. 997331 hat den quadratischen Rest $-53$, also müsste nach dem
Reziprozitätsgesetz $-997331$ ein Rest mod 53 sein. $-997331$ hat mod 53 den
Rest $-30$, $-30$ müsste daher quadratischer Rest mod 53, also $53-30=23$
quadratischer Rest mod 53 sein. Die Tafel zeigt aber 23 als Nichtrest, folglich kann 997331 keine Primzahl sein.

In dem zitierten Briefe an \textsc{Zimmermann} spricht \textsc{Gauss} auch davon, dass
man die \textls[120]{quadratischen Reste auf Stäbchen anbringen} könne, wodurch
die praktische Verwendbarkeit noch erhöht würde. Auch in den \textit{Disquisitiones}\footnote{Werke I, S. 404, 405.}
äussert er sich darüber: »Ad summum autem commoditatis fastigium usus talis
tabulae evehetur, si singulae columnellae verticales, e quibus constat, exsecantur
lamellisque aut baculis (\textsc{Neper}ianis similibus) agglutinantur, ita ut eae, quae
in quovis casu sunt necessariae, i. e. quae numeris $r,\;r',\;r''$ etc. residuis numeri propositi in factores resolvendi, respondent, separate examinari possint..«.
\textsc{Gauss} ist nun nicht bei der theoretischen Erörterung stehen geblieben, sondern er hat den Gedanken in die Tat umgesetzt und sich \textls[120]{einen solchen
Apparat selbst hergestellt}, der sich jetzt im \textsc{Gauss}archiv befindet. Wir
sind daher in der Lage, ihn auf Grund eigener Anschauung beschreiben zu
können.

Von den 13 gleichlangen, schmalen, mit Papier überklebten Holzstäbchen
ist das erste äusserst sorgfältig in gleiche Teile eingeteilt, so dass die Primzahlen von 7 bis 553 in gleichen Abständen darauf angebracht werden konnten.
Auf den übrigen 12 sind Streifen, wo ein quadratischer Rest ist, und leere
Stellen da, wo ein Nichtrest ist.

Merkwürdiger Weise sind es aber nicht Stäbchen für die aufeinanderfolgenden Primzahlen, sondern für $-1,\;2,\;-2,\;3,\;-3,\;5,\;-5,\;-6,\;7,\;11,$
$13,\;17$. Es scheint daraus hervorzugehen, dass \textsc{Gauss} einen Kunstgriff besass,
um einige von diesen Resten besonders bequem zu erhalten. Äusserungen
über einen derartigen Kunstgriff liessen sich nicht ermitteln, doch ist es höchst
wahrscheinlich, dass er die Fertigkeit besass, gegebene grosse Zahlen sehr
\clearpage\SourceLeaf{590}\ArithmeticHead{35}{R}{VII. DIE TAFEL DES QUADRATISCHEN CHARAKTERS DER PRIMZAHLEN.}
\setcounter{footnote}{0}
\noindent rasch in die Form $a^2+b^2$, sowie $a^2\pm xb^2\ (x=2,\;3,\;5,\;7,\;11,\;13,\;17)$ und endlich in die Form $2a^2+3b^2$ zu bringen. Fand er z. B. $n=a^2+5b^2$, so
multiplizierte er mit 5 und erhielt $5n=5a^2+25b^2$, $5n-5a^2=(5b)^2$, also
ist $-5a^2$, mithin auch $-5$ quadratischer Rest. Diese Vermutung gewinnt
dadurch an Wahrscheinlichkeit, dass \textsc{Gauss} dem allgemeinen Problem (\textit{Disqu.
Arithm.}, art. 146\textsuperscript{1}): »Propositis duobus numeris quibuscunque $P,\;Q$, invenire,
utrum alter $Q$, alterius $P$ residuum sit an non residuum« die Betrachtung von
Einzelfällen vorausschickt, und zwar: Residuum $-1$ (S. 82), Residua $+2$ et
$-2$ (S. 84), Residua $+3$ et $-3$ (S. 88), Residua $+5$ et $-5$ (S. 90), De $\pm7$
(S. 93).
\setcounter{footnote}{1}\footnotetext{Werke I, S. 111.}

Die Anfertigung der Stäbchen ist mit der wunderbaren Sorgfalt ausgeführt, die auch \textsc{Gauss}’ übrige Arbeiten kennzeichnet; die Striche und Punktierungen sind mit solcher Genauigkeit gemalt, dass man sie auf den ersten
Anblick für gedruckt halten möchte.

Auch die ausführliche Tafel, die \textsc{Schering} unter dem Titel \textit{Tafel des
quadratischen Charakters der Primzahlen}\footnote{Werke II, S. 399.} veröffentlicht hat (in der Handschrift
lautet der Titel: \textit{Quadratorum numeris primis divisorum residua lateralia}), ist
peinlich sorgfältig ausgeführt. \textsc{Schering} hat in der \textsc{Gauss}schen Tafel 190
Fehler gefunden, die er bei der Drucklegung berichtigte. Unsere Hoffnung,
aus der Verteilung dieser Fehler einen Schluss auf die bei der Berechnung
der Tafel angewandten Methoden ziehen zu können, hat sich nicht erfüllt,
nur eine bemerkenswerte Tatsache hat sich ergeben, nämlich die, dass \textsc{Gauss}
das \textls[120]{Reziprozitätsgesetz nicht zur Kontrolle benutzt} hat.

Es ist daher nicht unwahrscheinlich, dass er dieses Gesetz noch nicht
kannte, als er die Tabelle anfertigte, und dass er es erst während der Herstellung derselben induktiv fand. Allerdings ist auch die Möglichkeit nicht
abzuweisen, dass er hier überhaupt keine Kontrollrechnung vorgenommen hat.
Diese Frage wird uns noch in einem späteren Abschnitt beschäftigen.
\par\hfill 5\textsuperscript{*}
\clearpage\endgroup
